\chapter[Introdução]{Introdução}\label{capitulo1}

A avaliação psicopedagógica é um dos principais instrumentos de trabalho do psicopedagogo e tem como finalidade compreender as causas e as manifestações das dificuldades e dos transtornos de aprendizagem. Ela não se limita a identificar lacunas no desempenho escolar: procura entender como fatores cognitivos, emocionais, sociais e contextuais interferem no processo de aprender. É a partir dessa leitura que o profissional formula hipóteses diagnósticas e planeja intervenções adequadas a cada aprendente \cite{Santos_Silva_2020}. Por isso, o processo não se reduz à aplicação de testes padronizados, mas depende da articulação entre instrumentos, observações e a história do sujeito \cite{tavares2022_instrumentos_psicopedagogicos}.

No acompanhamento do desenvolvimento infantil, instrumentos como o Denver II e o \textit{Ages and Stages Questionnaires} (ASQ) são amplamente utilizados no rastreio de marcos do desenvolvimento e de possíveis sinais de atraso, inclusive em creches e pré-escolas \cite{Rezende_Costa_Pontes_2005, Ramos_DellaBarba_2021}. A escolha de instrumentos validados e a experiência de quem os aplica influenciam diretamente a qualidade das interpretações \cite{recima2024_instrumentos}. Outro aspecto relevante é a continuidade do registro: como a avaliação psicopedagógica também acompanha a criança ao longo do tempo, os registros precisam permitir comparar sua evolução entre os atendimentos \cite{Santos_Silva_2020}.

Apesar do repertório consolidado de métodos e protocolos, a organização digital desse processo ainda é limitada. É comum que os profissionais trabalhem com registros em papel, planilhas isoladas e documentos dispersos, o que dificulta a padronização, reduz a rastreabilidade das informações e prejudica a continuidade do acompanhamento \cite{anpae2024_avaliacao_educacao_infantil}. Essa foi também a situação relatada pela psicopedagoga entrevistada na etapa de levantamento de requisitos (Seção~\ref{sec:coleta_dados}). Há, portanto, espaço para soluções que tratem o próprio instrumento avaliativo de forma estruturada, permitindo que o profissional o adapte sem perder a referência das aplicações já realizadas.

Com a intensificação do uso de tecnologias digitais no período pós-pandemia, surgiram novos desafios e possibilidades para a educação. Na Educação Infantil, a incorporação desses recursos precisa respeitar o caráter lúdico, interativo e relacional das aprendizagens, bem como os princípios éticos que orientam o trabalho com crianças \cite{Domingues_Kresciglova_Paschoal_Furlan_2023}. Quando bem planejadas, plataformas digitais podem facilitar o registro das informações, a comunicação entre escola, família e profissionais e o acesso a um histórico organizado, desde que não descaracterizem o processo avaliativo psicopedagógico.

Na engenharia de software aplicada à educação, interfaces centradas no usuário e fluxos de dados bem definidos favorecem a organização das informações e a análise longitudinal dos registros \cite{Lemes_Dias_Oliveira_2023}. No contexto clínico-pedagógico, porém, a elaboração de relatórios interpretativos exige liberdade de escrita e contextualização, o que raramente é contemplado por sistemas genéricos de gestão escolar \cite{junior2024_gestao_escolar}. Uma ferramenta voltada à psicopedagogia precisa, então, equilibrar padronização e flexibilidade: estruturar o que pode ser estruturado sem engessar a análise qualitativa do profissional.

\citeonline{lemos2023_psicopedagogia_tecnologias_digitais} aponta como dificuldade recorrente na psicopedagogia a falta de ferramentas digitais que reúnam, em um mesmo ambiente, dados quantitativos e qualitativos, como resultados de testes, observações, entrevistas, devolutivas e relatórios de acompanhamento. Além disso, estudos sobre dificuldades específicas de aprendizagem, como a discalculia, mostram que os protocolos frequentemente precisam ser adaptados às particularidades de cada criança \cite{tavares2022_instrumentos_psicopedagogicos}, o que torna ainda mais importante registrar essas adaptações.

Diante desse cenário, este trabalho propõe o desenvolvimento de um sistema web de apoio à avaliação psicopedagógica na Educação Infantil que reúna três frentes: o cadastro de aprendentes, com acesso restrito ao profissional responsável; o cadastro de instrumentos, com possibilidade de adaptação e controle de versões; e a elaboração de relatórios interpretativos. Busca-se melhorar a organização, a rastreabilidade e a continuidade do processo avaliativo, sem substituir a análise do profissional.

O trabalho se justifica por três razões principais:

\begin{itemize}
    \item Há no mercado brasileiro sistemas de prontuário eletrônico voltados a psicopedagogos, psicólogos e clínicas multiprofissionais, mas o levantamento apresentado na Seção~\ref{sec:solucoes_mercado} mostra que eles priorizam a gestão do consultório e oferecem, no máximo, um conjunto fixo de instrumentos ou modelos genéricos de formulário. Nenhum deles informa permitir que o profissional cadastre e adapte seus próprios instrumentos com um controle de versões que preserve, em cada aplicação, a versão efetivamente utilizada.
    \item Reunir as avaliações em um único ambiente reduz a perda de dados e a fragmentação dos documentos, além de permitir acompanhar a evolução do aprendente e comparar diferentes momentos da avaliação.
    \item Como as conclusões psicopedagógicas são interpretativas e dependem do contexto, o profissional precisa de uma ferramenta que reúna dados quantitativos e qualitativos e que deixe espaço para relatórios individualizados.
\end{itemize}

\section{Objetivo Geral}

Desenvolver um sistema web multiusuário de apoio à avaliação psicopedagógica na Educação Infantil que permita o registro estruturado das informações dos aprendentes, o gerenciamento dos instrumentos avaliativos e o acompanhamento das etapas da avaliação, contribuindo para a padronização, a rastreabilidade e a qualidade das devolutivas elaboradas pelos profissionais.

\section{Objetivos Específicos}

\begin{enumerate}
    \item Identificar, por meio de revisão bibliográfica, os principais instrumentos e procedimentos utilizados na avaliação psicopedagógica infantil, considerando as dimensões cognitiva, emocional e pedagógica.
    \item Projetar a arquitetura do sistema, com módulos para a gestão de aprendentes, o cadastro e a aplicação de instrumentos avaliativos, o registro de observações qualitativas e a elaboração de relatórios.
    \item Implementar mecanismos de proteção das informações, como controle de acesso por perfil e por titularidade dos registros, armazenamento de senhas por \textit{hash} e registro de auditoria, e analisar a adequação do sistema à Lei Geral de Proteção de Dados (LGPD), identificando as medidas ainda necessárias para o uso com dados reais.
    \item Desenvolver uma interface orientada à usabilidade e à acessibilidade, adequada à rotina do profissional no contexto clínico-pedagógico.
    \item Avaliar o sistema com profissionais da psicopedagogia, verificando a adequação das funcionalidades às demandas do processo avaliativo.
    \item Discutir as contribuições e as limitações do sistema em relação aos trabalhos relacionados e à percepção das profissionais participantes.
\end{enumerate}

\section{Organização do Trabalho}

Além desta introdução, o trabalho está organizado em mais seis capítulos. O Capítulo~\ref{capitulo2} apresenta os estudos relacionados, com trabalhos sobre prontuários eletrônicos, triagem digital do desenvolvimento infantil, uso de \textit{software} na avaliação psicopedagógica e sistemas web aplicados à educação, além das soluções disponíveis no mercado. O Capítulo~\ref{capitulo3-0} reúne a fundamentação teórica sobre desenvolvimento infantil, avaliação psicopedagógica e seus instrumentos, e os conceitos de engenharia de software que orientaram as decisões técnicas. O Capítulo~\ref{capitulo4} descreve os materiais e métodos, incluindo a coleta de dados, as etapas de desenvolvimento e os procedimentos éticos. O Capítulo~\ref{capitulo5} detalha o desenvolvimento do sistema, dos requisitos à arquitetura, à implementação, à segurança e aos testes. O Capítulo~\ref{capitulo6} apresenta e discute os resultados, com o atendimento aos requisitos, a avaliação de acessibilidade e a avaliação com profissionais. Por fim, o Capítulo~\ref{capitulo7} traz as conclusões, as limitações e as sugestões de trabalhos futuros.